% NAF 5166, batch 2 — Gallica views 21–40.
% Pass: Opus 5.5 (claude-opus-5-5), 2026-09-23 — first pass, unchecked against the views by a human.
%
% What the views are. Scans of about 4707 × 6431–6463, portrait, one page per
% view, some in colour and some in greyscale. Each leaf is smaller than the
% frame and mounted on a guard, with the edges of the guards and the facing
% board showing at one side. Read from the local mirror only: a 1000-px
% thumbnail of every view, four full-width bands per written view at 2400 px
% (12 % overlap), and tight crops at 600–2400 px for every struck, blotted or
% doubtful stroke and for every page number. View 40 was also read rotated
% 180°. Nothing was fetched from Gallica. \page{N} is always the view.
%
% What the twenty views hold. One piece, in one hand, in French: the middle of
% a memoir on the attraction of a spheroid very little different from a
% sphere, and on Laplace's equation at its surface. The rectos are written, on
% one side of the leaf only; the versos are blank (views 22, 24 … 38: nothing
% but the recto seen through the paper and some ink spots). Every page runs
% straight on from the one before, usually mid-sentence.
%
%   v. 21     Page « 3 ». Begins with paragraph « 2. » — so the memoir
%             starts BEFORE view 21 (paragraph 1, with the first expression of
%             V the page refers to as « l'expression precedente de V », is on
%             earlier views, presumably in batch 1). The spheroid R = a(1+y');
%             V as the sphere's term 4πa³/3r plus u, a double integral over
%             π', θ'; « Je differentie maintenant comme M. Laplace ».
%   v. 23     Page « 4 ». du/dr at the surface; the equation
%             u/2 + adu/dr = 0, then V/2 + adV/dr = −2πa²/3 + 2πa²y against
%             « M. Laplace trouve … −2πa²/3 parce qu'il suppose que la sphere
%             touche le spheroide ». A six-line margin note, crossed out,
%             says the same. Paragraph 3 begins: y' constant.
%   v. 25     Page « 5 ». The case y' constant computed: u/2 + adu/dr =
%             −2πa²y', which contradicts the equation. Paragraph 4: ρ, the
%             inverse distance.
%   v. 27     Page « 6+ ». ρ/2 + adρ/dr = 0 « équation identique »; the
%             expansion of ρ in Laplace's descending series in R⁽⁰⁾, R⁽¹⁾ …
%   v. 29     Page « 6 bis ». The series R⁽⁰⁾ + 3R⁽¹⁾ + 5R⁽²⁾ + &c. vanishes
%             at r = a; paragraph 5: the integral of an identically null
%             function should vanish, « Cependant » —
%   v. 31     Page « 7 ». — the constant case says otherwise: « ce qui est un
%             paradoxe dans le Calcul integral ». Paragraph 6: the calculation
%             redone « sans rien negliger »; u/2 + rdu/dr =
%             −a³(r²−a²)/2 ∫ρ³y'…, whose second term does not always vanish
%             at r = a.
%   v. 33     Page « 8 ». Paragraph 7: y' constant, ∫ρ³… = 4πy'/r(r²−a²),
%             hence u/2 + adu/dr = −2a²πy' « comme on l'a trouvée dans le
%             n°. 3 ».
%   v. 35     Page « 9 ». Paragraph 8: the general case, the pole moved to
%             the line r, new angles φ, ψ, the spherical triangle —
%   v. 37     Page « 10 ». — its formulas; d'Alembert named (« C'est ainsi
%             que d'Alembert a calculé l'attraction des spheroides peu
%             differens de la sphere »); paragraph 9: integration by parts in
%             μ, with Y the value of y' at μ = −1.
%   v. 39     Page « 11 ». The general case gives again u/2 + adu/dr =
%             −2a²πy « comme dans le cas où y est constant »; the last line,
%             « Il est bon de remarquer qu'une nouvelle integration par
%             parties … n'ajouteroit rien », stops at the foot of the text
%             with no punctuation.
%   v. 40     Back of the leaf of v. 39: upside down, three lines struck with
%             four long diagonals, headed « 11 » (itself struck) — a first,
%             abandoned start of page 11, the same words as the head of
%             v. 39 — and a small pen sketch of two nested arcs. Transcribed
%             as struck text.
%
%   Continuity at the edges of the batch.
%   (a) View 21 CONTINUES a text begun before it: it is page 3 and opens on
%       paragraph 2, referring back to an expression of V given earlier.
%   (b) Whether view 40 runs on into view 41 cannot be told from this
%       batch. View 40 is a cancelled draft on the back of page 11 and ends
%       on a struck « ainsi »; the live text ends on view 39 with
%       « … n'ajouteroit rien », unpunctuated, at the foot of the writing
%       with the rest of the page blank. If the memoir goes on, its next
%       page should carry « 12 » and open with paragraph 9's continuation
%       or a paragraph 10; the coordinator should look at the head of v. 41.
%
%   Who. The leaves name « M. Laplace » (v. 21, 23, 27) and d'Alembert
%   (v. 37), and speak in the first person (« je vais reprendre », « je
%   ferai »). No signature, no title, no date in this batch. The catalogue
%   gives the volume as Lagrange's papers; nothing in these twenty views
%   names the writer, and the pass attributes nothing.
%
%   The hand. One regular, rapid eighteenth-century cursive throughout, in
%   dark ink, the same pen for text and page numbers; long terminal strokes;
%   deletions made by a heavy line, by solid blots or by close vertical
%   hatching (v. 23, 31, 33). Accents are sparse (« spheroide »,
%   « differentie », « integrale » unaccented; « très », « à » accented);
%   the pass gives an accent where the band showed one, and at band scale it
%   may have missed or supplied a few: a check against the leaf is wanted.
%   Spelling kept: « complette », « parceque », « ensorte », « devroit »,
%   « disparoit », « disparoisse », « n'ajouteroit », « differens »,
%   « coté », « egale » for the radius.
%
%   Notation, held to throughout.
%   – The differential sign is the writer's ordinary looped d, the same
%     glyph as in « dans », « devient » (compared on v. 23), for total and
%     partial differentials alike; set d. This pass first set a round ∂; the
%     coordinator changed it to d, as batches 1 and 3 read the same hand.
%   – THE HAND'S SMALL 2 IS FORMED LIKE AN « n » (or like its r) with no
%     loop: in exponents (r², a²) and as the denominator in u/2, V/2, ρ/2,
%     −a/2√2. Batch-level evidence: V/2 + adV/dr = 2πa³/3r − 4πa⁴/3r²
%     (v. 23) and u/2 + adu/dr = 2πa²y' − 4πa²y' = −2πa²y' (v. 25) only
%     hold with 2; on v. 31 and 39 the same 2 is written fully.
%   – π is both 180° (« en nommant π l'angle de 180° », v. 21) and the
%     longitude of the attracted point; π' and θ' the angles of the
%     attracting molecule, θ the attracted point's. θ is a slanting oval
%     without stem or tail; never read β.
%   – μ is written with a long tail, as μ = cos θ' after the pole is moved.
%   – φ and ψ (from v. 35): a loop crossed by a stem, and a cross with a
%     hooked right arm. The capital Y (v. 37, 39) is sometimes a barred Y
%     and sometimes exactly the hooked cross of ψ; they are told apart by
%     the sense only, and the notes say where.
%   – The radical is a large « V » or « U » before a bracket; kept as
%     \sqrt{( … )}, including where the bracket also carries the exponent
%     3/2 (v. 23), and with its bar as short as the leaf draws it (v. 27).
%   – « &c. » is written as an α with two dots; set \&c.
%   – Signs as written: several minus signs the calculation needs are not on
%     the leaf (du/dr at v. 23, dρ/dr at v. 27); noted, not supplied.
%
% Numbers on the leaves. One series, top right of each written page, in
%   ink: « 3 » (v. 21), « 4 » (v. 23), « 5 » (v. 25), « 6+ » (v. 27),
%   « 6 bis » (v. 29), « 7 » (v. 31), « 8 » (v. 33), « 9 » (v. 35),
%   « 10 » (v. 37), « 11 » (v. 39), and a struck « 11 » heading the
%   cancelled draft on the back of that leaf (v. 40).
%   This pass first took the series for the writer's pagination, since it
%   starts at 3 and repeats 6. Read across the whole volume it is the
%   library's count of leaves: batch 1 has « A. », « 1 », « 2 » (v. 15–19),
%   batch 3 has 12 to 22 (v. 41–60) across four different hands, batch 4
%   has 23 to 26 (v. 62–68), and the certificate of 1888 (v. 13) announces
%   « 26 Feuillets plus les Feuillets A préliminaire & 6 bis » — the « 6 bis »
%   of v. 29. Because it is in ink, as in NAF 5161, no \folio{} is written
%   (coordinator's reconciliation, 2026-09-23). The « + » after the 6 on
%   v. 27 remains unexplained.
%   Other marks: a small ink sign like a « c » under the 9 (v. 35); a
%   boxed « I 65 » in a thinner pen in the right margin of v. 37, with an
%   opening bracket « [ » before « Ainsi la formule » on the same line —
%   meaning unknown (a later reader's reference? a copying mark?), recorded
%   as a \marginal{} with a note; paragraph numbers 2 (v. 21), 3 (v. 23),
%   4 (v. 25), 5 (v. 29), 6 (v. 31), 7 (v. 33), 8 (v. 35), 9 (v. 37),
%   some with a point, some without; references « n°. 3 » (v. 33, 35),
%   « n°. 4 » (v. 37), « n°. 6 » (v. 39) are to these paragraphs.
%
% Views not transcribed: 22, 24, 26, 28, 30, 32, 34, 36, 38 — blank versos
% (the recto seen through, mirrored and checked on v. 32; ink spots and
% bleed-through only). View 40 is transcribed although it is a verso,
% because it is written.

\documentclass[11pt,a4paper]{article}
% The preamble shared by every transcription of the mathematicians' manuscripts in Gallica.
%
% It rests on one idea: **uncertainty compounds, it does not resolve.** A
% transcription that smooths over a doubtful word has destroyed the only thing
% it was made for — a reader must be able to tell, at every point, what was on
% the page and what was supplied. The apparatus macros below are therefore the
% critical apparatus, not decoration.
%
% The same commands are recognised by `scripts/render.mjs`, which produces the
% site's reading view, and by `scripts/tei.mjs`. A macro added here must be
% added there too, or rendering fails loudly — which is the wanted behaviour.
%
% Comments here are English, like the rest of the repository. The documents
% this preamble sets are French, and that is deliberate: see the skills.

\ProvidesPackage{germain}[2026/09/15 Transcriptions of mathematicians' manuscripts in Gallica]

\usepackage{amsmath,amssymb,amsthm}
\usepackage{mathrsfs}
% Kept from the parent preamble so the renderer's diagram support stays
% available; these manuscripts draw no commutative diagrams, and nothing here uses it.
\usepackage{tikz-cd}
\usepackage[english,french]{babel}
\usepackage{fontspec}

% --- Greek ------------------------------------------------------------------
% The text face stays fontspec's default, Latin Modern, which has no Greek:
% XeTeX drops every glyph it lacks with a line in the log and nothing on the
% page, and the Greek words the manuscripts quote vanished from the PDFs. So
% the Greek and Greek Extended blocks get a XeTeX character class of their own,
% and entering or leaving a run of it switches the family to CMU Serif — the
% same Computer Modern design, with polytonic Greek — and back. Only the
% family changes: a Greek word in \textit or \textbf keeps its shape and series.
% The transitions to and from every other class carry the tokens that class
% has towards an ordinary letter, so French spacing before « ; » or inside
% « » still holds after a Greek word. No grouping, so no brace or \endgroup
% can come out unbalanced; maths is left alone, since XeTeX inserts nothing
% there. Kept identical in `exercices.sty`.
\makeatletter
\newfontfamily\germainGreekfont{cmun}[
  Extension=.otf, NFSSFamily=germaingreek,
  UprightFont=*rm, ItalicFont=*ti, SlantedFont=*sl,
  BoldFont=*bx, BoldItalicFont=*bi, BoldSlantedFont=*bl]
\newXeTeXintercharclass\germain@greek
\def\germain@greekrange#1#2{%
  \@tempcnta=#1\relax
  \loop
    \XeTeXcharclass\@tempcnta=\germain@greek
    \ifnum\@tempcnta<#2\advance\@tempcnta\@ne\repeat}
\germain@greekrange{"0370}{"03FF}
\germain@greekrange{"1F00}{"1FFF}
\def\germain@greekprev{\rmdefault}
\def\germain@greekin{%
  \edef\germain@greekprev{\f@family}\fontfamily{germaingreek}\selectfont}
\def\germain@greekout{\fontfamily{\germain@greekprev}\selectfont}
\def\germain@greekjoin{%
  \XeTeXinterchartokenstate=\@ne
  \@tempcnta=\z@
  \loop
    \ifnum\@tempcnta=\germain@greek\else
      \XeTeXinterchartoks\germain@greek\@tempcnta=\expandafter{%
        \expandafter\germain@greekout\the\XeTeXinterchartoks\z@\@tempcnta}%
      \XeTeXinterchartoks\@tempcnta\germain@greek=\expandafter{%
        \the\XeTeXinterchartoks\@tempcnta\z@\germain@greekin}%
    \fi
    \ifnum\@tempcnta<4095\advance\@tempcnta\@ne\repeat}
% babel-french sets its own classes, and saves and restores the interchar
% state, each time a language is selected: join again after it does.
\addto\extrasfrench{\germain@greekjoin}
\addto\extrasenglish{\germain@greekjoin}
\AtBeginDocument{\germain@greekjoin}
\makeatother

\usepackage{geometry}
\geometry{a4paper,margin=25mm,marginparwidth=28mm}
\usepackage{marginnote}
\usepackage{xcolor}
\usepackage[normalem]{ulem}
\setlength{\emergencystretch}{3em}
\usepackage{hyperref}
\hypersetup{colorlinks=true,linkcolor=black,urlcolor=black,pdfborder={0 0 0}}

% --- Batch metadata ---------------------------------------------------------
% Taken as they stand by the rendering script, which shows them at the head of
% the reading view. They fix what the file claims to cover. The macro names are
% the parent project's, so that the scripts stay shared; \folder here means the
% volume's slug — `fr-9115` — and \pages counts Gallica views.

\newcommand{\folder}[1]{\gdef\grothFolder{#1}}
\newcommand{\batch}[1]{\gdef\grothBatch{#1}}
\newcommand{\pages}[2]{\gdef\grothFirst{#1}\gdef\grothLast{#2}}
\newcommand{\foldertitle}[1]{\gdef\grothFoldertitle{#1}}

% The holder's own shelfmark, printed as they write it: « Français 9115 ».
\newcommand{\shelfmark}[1]{\gdef\grothShelfmark{#1}}

% The dating, copied verbatim from the holder's catalogue — never inferred
% here. « XIXe siècle » is the BnF's; a pass that dates a leaf from its content
% says so in a \note{}, as its own inference.
\newcommand{\dating}[1]{\gdef\grothDating{#1}}

% The status notice, printed in the title block and repeated as a diagonal
% watermark on every page of the PDF. The manuscripts are in the public
% domain; what this line declares is not a right but a quality — an
% unverified first pass by a machine. The skills write the line; a file
% without it gets no watermark, so leaving it out is visible in the output.
\newcommand{\watermark}[1]{\gdef\grothWatermark{#1}}

\gdef\grothFolder{?}\gdef\grothBatch{}\gdef\grothFirst{?}\gdef\grothLast{?}%
\gdef\grothFoldertitle{}\gdef\grothShelfmark{}\gdef\grothDating{}\gdef\grothWatermark{}

% --- The résumé -------------------------------------------------------------
\newenvironment{resume}
  {\small\setlength{\parskip}{0.45em}}
  {\par\medskip\hrule\medskip}

% --- Keywords ---------------------------------------------------------------
% In English, closing the modernised reading's résumé: the modern vocabulary
% under which to search for what the leaves construct. The manifest extracts
% them to tag the volume — this line is the single source of the tags.
\newcommand{\keywords}[1]{\par\smallskip\noindent
  {\footnotesize\textsc{Keywords} --- \textit{#1}}\par}

% --- The critical apparatus -------------------------------------------------

% The Gallica view number, in the margin. This is the anchor that lets a
% disagreement over a formula be settled by looking at *one* image rather than
% a volume of 750. The site uses it to turn the facsimile as the transcription
% is scrolled. A view is one scanned image: usually one side of a leaf.
\newcommand{\page}[1]{%
  \marginnote{\footnotesize\color{gray!70}\textsf{v.\,#1}}%
  \label{page:#1}\ignorespaces}

% The BnF's pencilled foliation, where the leaf carries it and a pass can read
% it: « 198r ». This is what the literature cites — « ff. 198r–208v » — and
% the one line that turns a citation into a view. Recorded, never inferred:
% a folio a pass cannot read is not written.
\newcommand{\folio}[1]{%
  \marginnote{\footnotesize\color{gray!50}\textsf{f.\,#1}}\ignorespaces}

% The modernised reading's coarser counterpart to \page{}: which views a
% section is drawn from.
\newcommand{\pagerange}[2]{%
  \marginnote{\footnotesize\color{gray!70}\textsf{v.\,#1--#2}}%
  \ignorespaces}

% Illegible: never guessed.
\newcommand{\ill}{\textcolor{red!60!black}{[\,\ldots\,]}}

% A doubtful reading: the word is given, and flagged as uncertain.
\newcommand{\uncertain}[1]{\uline{#1}}

% An editorial addition — a missing letter, an implied word.
\newcommand{\add}[1]{\textcolor{blue!50!black}{[#1]}}

% Struck out by the writer. What was crossed out is part of the document.
\newcommand{\struck}[1]{\sout{#1}}

% The transcriber's note, on the state of the page or the mathematics.
\newcommand{\note}[1]{\par\smallskip{\small\color{gray!60!black}\textit{[#1]}}\par\smallskip}

% A marginal note of hers, distinct from the body of the text.
\newcommand{\marginal}[1]{\par\smallskip\noindent\hspace{1em}%
  {\small\color{gray!60!black}#1}\par\smallskip}

% --- Title ------------------------------------------------------------------
% The title block is French, like the documents themselves.

\AddToHook{shipout/background}{%
  \ifx\grothWatermark\empty\else
    \begin{tikzpicture}[remember picture, overlay]
      \node[rotate=52, scale=3.4, align=center, text=gray!18,
            font=\sffamily\bfseries]
        at (current page.center) {\grothWatermark};
    \end{tikzpicture}%
  \fi}

\AtBeginDocument{%
  \begin{center}
    {\large\bfseries Manuscrits de mathématiciens (Gallica) ---
      \ifx\grothShelfmark\empty\grothFolder\else\grothShelfmark\fi}\\[2pt]
    {\itshape\grothFoldertitle}\\[4pt]
    {\small vues \grothFirst--\grothLast
      \ifx\grothBatch\empty\else\ (lot \grothBatch)\fi}\\[2pt]
    \ifx\grothDating\empty\else
      {\small\color{gray!60!black}Datation du catalogue : \grothDating}\\[2pt]
    \fi
    \ifx\grothWatermark\empty\else
      {\small\bfseries\color{red!45!gray}\grothWatermark}\\[2pt]
    \fi
    {\footnotesize\color{gray!60!black}Transcription établie d'après les images IIIF de Gallica.
     Source gallica.bnf.fr / Bibliothèque nationale de France.\\
     Les lectures incertaines sont soulignées, les illisibles notées [\,\ldots\,],
     les ajouts entre crochets ; \textsf{v.} numérote les vues de Gallica,
     \textsf{f.} la foliotation au crayon de la BnF.}
  \end{center}
  \vspace{1em}}

\endinput


\folder{naf-5166}
\shelfmark{NAF 5166}
\batch{2}
\pages{21}{40}
\dating{1701-1900}
\watermark{Lecture automatique\\non vérifiée}
\foldertitle{Papiers du géomètre J.-L. LAGRANGE (1813).}

\begin{document}

\page{21}
\note{Recto d'un feuillet monté sur onglet. En haut à droite, un grand « 3 » à l'encre noire, du même trait appuyé que le texte : voir l'en-tête pour ce que peut être cette série. La page commence par le paragraphe numéroté « 2. » : le paragraphe 1 et la première expression de $V$, à laquelle renvoie « l'expression precedente de $V$ », sont avant cette vue.}

2. Supposons maintenant que le spheroide soit très peu different d'une sphere dont le rayon soit egale à $a$, on aura $R = a(1+y')$, $y'$ étant une fonction très petite des sinus et cosinus de $\pi'$ et $\theta'$, dont on negligera le carré et les puissances superieures; il est clair que la valeur de $V$ sera egale a sa valeur relativement a la sphere, plus à la valeur relative a l'excès du spheroide sur la sphere. On sait et il est facile de trouver que la premiere est egale a la masse de la sphere divisée par la distance $r$; elle est ainsi \struck{\uncertain{egale}} exprimée par $\frac{4\pi a^3}{3r}$, en nommant $\pi$ l'angle de $180^\circ$. Quant à la seconde, comme nous ne tenons compte que des premieres dimensions de $y$, il est aisé de voir, qu'on l'aura en mettant dans l'expression precedente de $V$, $a$ à la place de $R$ et $ay$ à celle de $dR$. Ainsi en nommant cette seconde partie $u$, on aura
\[ V = \frac{4\pi a^3}{3r} + u, \quad\text{et}\quad u = \struck{\ill{}}\, a^3 \int \frac{y'\,d\pi'\,d\theta' \sin\theta'}{\sqrt{(r^2 - 2ra\mu + a^2)}} \]
où il n'y a plus que deux integrations a faire l'une relative a $\pi'$ l'autre relative a $\theta'$ ou à $\mu$.
\note{« 2. » : le chiffre est fermé d'une boucle à sa base, lu 2. « egale » (le rayon) : ainsi sur la page. Le mot biffé après « ainsi » est rayé d'un trait épais ; on y devine « egale », répété de la ligne précédente. « $\theta'$ » : une boucle ovale inclinée, sans hampe ni jambage, ici et dans tout le lot ; la lecture $\beta'$ est exclue par l'absence de jambage. La lettre qui suit $d\pi'$ au numérateur est récrite sur une autre. Entre « $u =$ » et $a^3$, un signe ou un coefficient noyé sous une grosse tache d'encre, effacement de l'auteur. « $ay$ » : sans prime sur l'$y$, comme « des premieres dimensions de $y$ ». Le signe de différentielle est partout un $d$ rond, gardé.}

Je differentie maintenant comme M. Laplace par rapport à $r$, j'ai
\[ \frac{d V}{d r} = -\frac{4\pi a^3}{3r^2} + \frac{d u}{d r} \]
\note{La page s'arrête sur cette équation ; la suite est en tête de la vue 23 (« Or $\frac{d u}{d r} = \ldots$ »).}

\page{23}
\note{En haut à droite, « 4 » à l'encre, du trait de la vue 21. La page continue sans rupture l'équation qui termine la vue 21.}

Or
\[ \frac{d u}{d r} = a^3 \int \frac{(r - a\mu)\,y'\,d\pi'\,d\mu}{\sqrt{(r^2 - 2ar\mu + a^2)}^{\,\frac{3}{2}}}. \]
Lorsque le point attiré est à la surface du spheroide on a $r = a(1+y)$, en designant par $y$ ce que devient $y'$ lorsque $\theta'$ et $\pi'$ deviennent $\theta$ et $\pi$; mais comme nous negligeons dans $u$ les quantités du second ordre, il suffit de faire $r = a$. On aura ainsi:
\[ u = \frac{a^2}{\sqrt{2}} \int \frac{y'\,d\pi'\,d\mu}{\sqrt{(1-\mu)}}, \quad \frac{d u}{d r} = \frac{-a}{2\sqrt{2}} \int \frac{(1-\mu)\,y'\,d\pi'\,d\mu}{(1-\mu)^{\frac{3}{2}}} = -\frac{a}{2\sqrt{2}} \int \frac{y'\,d\pi'\,d\mu}{\sqrt{(1-\mu)}}; \]
\note{Première équation : pas de signe moins devant $a^3$, ainsi sur la page ; le dénominateur porte le même signe radical que partout ailleurs, suivi d'une parenthèse fermée par un crochet et de l'exposant $\frac{3}{2}$ : transcrit tel quel. Sous $u$, $V$ et $-a$, le dénominateur est le petit 2 de cette main, fait comme un n sans boucle, qu'on prendrait pour $r$ ; c'est le 2 des exposants ($r^2$, $a^2$), et le calcul le confirme ($\frac{V}{2} = \frac{2\pi a^3}{3r}$ plus bas). Le « 1 » de $(1-\mu)^{\frac32}$ est empâté. Après le deuxième membre, le signe d'égalité est surchargé.}

Donc
\[ \frac{u}{2} + \frac{a\,d u}{d r} = 0 \]
equation qui doit avoir lieu à la surface du spheroide.

Donc on aura aussi à cette surface \struck{\ill{}}
\[ \frac{V}{2} + \frac{a\,d V}{d r} = \frac{2\pi a^3}{3r} - \frac{4\pi a^4}{3r^2} \]
où il faut faire $r = a(1+y)$ ce qui donnera
\[ \frac{V}{2} + \frac{a\,d V}{d r} = -\frac{2\pi a^2}{3} + 2\pi a^2 y. \]
\marginal{\struck{\uncertain{avoit} pas pour \ill{} \uncertain{remplir} necessairement la condition que la sphere touche la surface du spheroide.}}
\note{Après « cette surface », un petit groupe de lettres noyé d'encre, biffé. Dans la marge de droite, en regard de ces deux équations, une note de six lignes, de la main du texte, barrée de trois grands traits croisés ; sa première ligne commence au bord (« avoit », peut-être précédé d'un « n' » caché) et le deuxième mot de la deuxième ligne n'a pas été lu. Elle reprend l'idée de la phrase sur M. Laplace qui suit.}

M. Laplace trouve \struck{\ill{}}
\[ \frac{V}{2} + \frac{a\,d V}{d r} = -\frac{2\pi a^2}{3} \]
parce qu'il suppose que la sphere touche le spheroide, auquel cas $y = 0$. \struck{\uncertain{Cependant comme} \ill{}}
\marginal{\struck{\ill{}}}
\note{Avant $\frac{V}{2}$, une petite tache biffée. La fin de la ligne après « $y = 0$. » est rayée d'un trait continu ; on y lit « Cependant comme » avec doute, et un dernier mot noyé. Dans la marge de droite, à la même hauteur, un mot entièrement noyé sous un trait épais.}

3. Considerons l'equation
\[ \frac{u}{2} + \frac{a\,d u}{d r} = 0, \]
qui est independante de la valeur de $y'$ dans $u$; et supposons d'abord $y'$ constant, on aura
\[ u = a^3 y' \int \frac{d\pi'\,d\theta' \sin\theta'}{\sqrt{(r^2 - 2ra\mu + a^2)}}; \]
il est facile de
\note{Le signe d'égalité de la dernière formule est repassé. Le $\theta'$ du numérateur est récrit sur une autre lettre, comme à la vue 21. La page s'arrête sur « il est facile de », au milieu de la phrase, le dernier mot à demi barré par le trait qui le prolonge ; le bas du feuillet est blanc. La suite est en tête de la vue 25.}

\page{25}
\note{En haut à droite, « 5 » à l'encre. La page reprend la phrase laissée à la fin de la vue 23 (« il est facile de »).}

d'avoir l'integrale de $\frac{d\pi'\,d\theta' \sin\theta'}{\sqrt{(r^2 - 2ra\mu + a^2)}}$, car en transportant l'axe, ou le pole des angles $\pi'$, $\theta'$ dans la ligne $r$, ce qui est permis, puisque cette ligne est fixe par rapport aux memes angles, on a $\theta = 0$, ce qui donne $\mu = \cos\theta'$, et reduit la differentielle dont il s'agit à
\[ -\frac{d\pi'\,d\mu}{\sqrt{(r^2 - 2ra\mu + a^2)}}, \]
laquelle doit être integrée depuis $\pi' = 0$ jusqu'a $\pi' = 2\pi$, et depuis $\mu = 1$ jusqu'a $\mu = -1$. La premiere integration donne
\[ -\frac{2\pi\,d\mu}{\sqrt{(r^2 - 2ra\mu + a^2)}}; \]
la seconde donne d'abord $\frac{2\pi\sqrt{(r^2 - 2ra\mu + a^2)}}{ra}$, et completant
\[ \frac{2\pi(r+a)}{ra} - \frac{2\pi(r-a)}{ra} = \frac{4\pi}{r}. \]
Ainsi on a
\[ u = \frac{4\pi a^3 y'}{r}, \quad\text{et de là}\quad \frac{d u}{d r} = -\frac{4\pi a^3 y'}{r^2}; \]
et faisant $r = a$, $u = 4\pi a^2 y'$, $\frac{d u}{d r} = -4\pi a y'$; mais ces valeurs ne satisfont pas à l'equation $\frac{u}{2} + \frac{a\,d u}{d r} = 0$; car elles donnent
\[ \frac{u}{2} + \frac{a\,d u}{d r} = -2\pi a^2 y'. \]
\note{Dans « $\frac{u}{2}$ », le 2 est le petit 2 en forme de n (voir la vue 23). La première fraction du numérateur, au début de la page, a le $\theta'$ récrit comme aux vues 21 et 23.}

4 Nommons $\rho$ pour abreger le radical $\frac{1}{\sqrt{(r^2 - 2ra\mu + a^2)}}$, ensorte que l'on ait $u = a^3 \int \rho\, y'\,d\pi'\,d\theta' \sin\theta'$; comme $r$ n'est contenu que dans $\rho$, on aura
\[ \frac{u}{2} + \frac{a\,d u}{d r} = a^3 \int \left(\frac{\rho}{2} + \frac{a\,d\rho}{d r}\right) y'\,d\pi'\,d\theta' \sin\theta', \]
\struck{\ill{}}en faisant $r = a$ après la differentiation.
\note{« 4 » : sans point, suivi d'un petit trait avant « Nommons ». « ensorte » : en un mot. La dernière ligne commence par une lettre noyée sous une tache d'encre, qui semble couvrir une initiale récrite en « en ». La phrase s'arrête là, aux trois quarts de la page ; le bas du feuillet est blanc. La vue 27 commence par « Or ».}

\page{27}
\note{En haut à droite, « 6+ » à l'encre : un 6, puis une petite croix. Le feuillet suivant (vue 29) porte « 6 bis ». La page continue le paragraphe 4 de la vue 25.}

Or
\[ \frac{d\rho}{d r} = \frac{r - a\mu}{(r^2 - 2ar\mu + a^2)^{\frac{3}{2}}}; \]
faisant $r = a$ on a
\[ \rho = \frac{1}{a\sqrt{2}(1-\mu)}, \quad \frac{d\rho}{d r} = \frac{1}{2a\sqrt{2}(1-\mu)}; \]
donc
\[ \frac{\rho}{2} + \frac{a\,d\rho}{d r} = 0 \]
équation identique.
\note{Les deux valeurs de $\frac{d\rho}{d r}$ sont écrites sans signe moins, ainsi sur la page, comme celle de $\frac{d u}{d r}$ à la vue 23 ; avec elles l'« équation identique » ne s'annulerait pas. La barre du radical ne couvre que le 2, et $(1-\mu)$ est écrit à sa suite : ainsi transcrit ; la valeur de $\rho$ demanderait $\sqrt{2(1-\mu)}$.}

Si on reduit, comme l'a fait M. Laplace, le radical $\rho$ en une serie descendante par rapport à $r$, on aura
\[ \rho = \frac{R^{(0)}}{r} + \frac{aR^{(1)}}{r^2} + \frac{a^2R^{(2)}}{r^3} + \&c. \]
\[ \frac{r\,d\rho}{d r} - \frac{R^{(0)}}{r} - \frac{2aR^{(1)}}{r^2} - \frac{3a^2R^{(2)}}{r^3} + \&c. \]
d'où resulte
\[ \frac{\rho}{2} + \frac{r\,d\rho}{d r} = -\frac{1}{2}\left(\frac{R^{(0)}}{r} + \frac{3aR^{(1)}}{r^2} + \frac{\uncertain{5}a^2R^{(2)}}{r^3} + \&c.\right) \]
et faisant $r = a$
\[ \frac{\rho}{2} + \frac{a\,d\rho}{d r} = -\frac{1}{2a}\left(R^{(0)} + 3R^{(1)} + 5R^{(2)} + \&c.\right) \]
serie dans laquelle tous les termes doivent se detruire d'eux memes; ce qu'on trouve en effet après la substitution des valeurs de $R^{(0)}$, $R^{(1)}$ \&c. en fonctions de $\mu$.
\note{« \&c. » : un signe en forme d'alpha suivi de deux points, l'abréviation de la main. Dans la deuxième ligne, entre $\frac{r\,d\rho}{d r}$ et $\frac{R^{(0)}}{r}$, un simple trait, sans signe d'égalité, et tous les termes précédés de « $-$ », tels qu'on les lit. Le 3 de « $3aR^{(1)}$ » est empâté ; le coefficient suivant, lu 5 avec doute, pourrait être un 3, mais la ligne d'après écrit nettement « $5R^{(2)}$ ». Les indices de $R$ : un 0 fermé et un 1, entre parenthèses.}

On n'a pas même besoin pour s'en convaincre d'effectuer ces substitutions; car puisque $\rho = (r^2 - 2ar\mu + a^2)^{-\frac{1}{2}}$, on aura $\rho\sqrt{r} = \left(r + \frac{1}{r} - 2a\mu\right)^{-\frac{1}{2}}$, et differentiant par rapport à $r$,
\[ \frac{\rho}{2\sqrt{r}} + \frac{d\rho}{d r}\sqrt{r} = -\frac{1}{2}\left(1 - \frac{a^2}{r^2}\right)\left(r + \frac{a^2}{r} - 2a\mu\right)^{-\frac{3}{2}}; \]
donc multipliant par $\sqrt{r}$
\note{« $r + \frac{1}{r}$ » : le numérateur est un petit trait vertical, lu 1 ; la ligne suivante écrit $\frac{a^2}{r}$ au même endroit. La page s'arrête sur « par $\sqrt{r}$ », sans ponctuation ; le bas du feuillet est blanc.}

\page{29}
\note{En haut à droite, « 6 bis » à l'encre, de la main du texte. La page reprend la phrase laissée à la fin de la vue 27 (« donc multipliant par $\sqrt{r}$ »).}

\[ \frac{\rho}{2} + \frac{r\,d\rho}{d r} = -\frac{r^2 - a^2}{2(r^2 - 2ar\mu + a^2)^{\frac{3}{2}}} = -\frac{(r^2 - a^2)\rho}{2(r^2 - 2ar\mu + a^2)} \]
de sorte qu'on aura identiquement
\[ \frac{R^{(0)}}{r} + \frac{3aR^{(1)}}{r^2} + \frac{5a^2R^{(2)}}{r^3} + \&c. = \frac{1 - \frac{a^2}{r^2}}{1 - \frac{2a\mu}{r} + \frac{a^2}{r^2}} \times \left(\frac{R^{(0)}}{r} + \frac{aR^{(1)}}{r^2} + \frac{a^2R^{(2)}}{r^3} + \&c.\right) \]
par consequent en faisant $r = a$ la serie
\[ R^{(0)} + 3R^{(1)} + 5R^{(2)} + \&c. \]
sera necessairement nulle.
\note{Les deux dénominateurs de la première ligne commencent par le petit 2 en forme de n. « $3R^{(1)}$ » : entre le 3 et le $R$, un trait d'attaque en forme de parenthèse ouvrante, que rien ne ferme ; non transcrit comme parenthèse.}

5. La fonction $\frac{\rho}{2} + \frac{a\,d\rho}{d r}$ de $\mu$ est donc toujours identiquement nulle. Donc l'integrale de cette fonction multipliée par $y'\,d\pi'\,d\theta' \sin\theta'$ étant supposée nulle lorsque $\pi' = 0$ et $\theta' = 0$ devroit être toujours nulle par les principes du calcul integral, suivant lesquels l'integrale est regardée comme la somme de toutes les valeurs successives de la differentielle. Ainsi on devroit avoir toujours l'equation $\frac{u}{2} + \frac{a\,d u}{d r} = 0$ que nous avons trouvée d'abord. Cependant
\note{La page s'arrête sur « Cependant », au milieu de la phrase, aux trois quarts du feuillet ; le bas est blanc. La suite est en tête de la vue 31.}

\page{31}
\note{En haut à droite, « 7 » à l'encre ; un peu à sa gauche et plus bas, une petite croix d'encre, comme celle qui suit le 6 à la vue 27. La page reprend la phrase laissée sur « Cependant » à la fin de la vue 29.}

venons de voir que dans le cas le plus simple où $y'$ est constant on a $\frac{u}{2} + \frac{a\,d u}{d r} = -2\pi a^2 y'$; ce qui est \struck{\ill{}} un paradoxe dans le Calcul integral.
\note{« venons » : sans sujet sur cette page ; le « nous » attendu après « Cependant » n'est ni au pied de la vue 29 ni en tête de celle-ci. Le mot biffé après « est » est court et couvert d'un trait épais.}

6. Pour en trouver l'explication je vais \struck{d'abord} reprendre l'analyse qui a donné l'equation $\frac{u}{2} + \frac{a\,d u}{d r} = 0$, et je ferai d'abord le calcul sans rien negliger. Puisque $u = a^3 \int \rho\, y'\,d\pi'\,d\theta' \sin\theta'$, et que la quantité $r$ n'est contenue que dans l'expression de $\rho$, il n'y a nul doute qu'en differentiant par rapport à $r$ on n'ait
\[ \frac{d u}{d r} = a^3 \int \frac{d\rho}{d r}\, y'\,d\pi'\,d\theta' \sin\theta'. \]
Or
\[ \frac{d\rho}{d r} = -\struck{\ill{}}\,\frac{r - a\mu}{(r^2 \uncertain{-} 2ar\mu + a^2)^{\frac{3}{2}}}; \]
donc
\[ \frac{r\,d\rho}{d r} = \frac{-\struck{\ill{}}\,1}{2\sqrt{(r^2 - 2ra\mu + a^2)}}\ \struck{\ill{}} - \struck{\ill{}}\,\frac{r^2 - a^2}{2(r^2 - 2ra\mu + a^2)^{\frac{3}{2}}} = \struck{\ill{}} -\frac{\rho}{2}\ \struck{\ill{}} - \frac{(r^2 - a^2)\rho^3}{2}; \]
donc on aura
\[ \frac{u}{2} + \frac{r\,d u}{d r} = -\struck{\ill{}}\,\frac{a^3(r^2 - a^2)}{2} \int \rho^3 y'\,d\pi'\,d\theta' \sin\theta'. \]
\note{« d'abord » après « je vais » est rayé d'un trait qui passe au pied des lettres ; il revient à la ligne suivante. Ces trois lignes de calcul sont très corrigées. Dans $\frac{d\rho}{d r}$, une tache biffée entre le signe moins et la fraction ; le signe entre $r^2$ et $2ar\mu$ est surchargé (un moins repassé en plus, ou l'inverse), lu moins avec doute. Dans $\frac{r\,d\rho}{d r}$ : au numérateur du premier terme, une tache d'encre avant le 1 ; après le premier terme, une grosse tache puis un signe moins ; une tache au début du numérateur $r^2 - a^2$ ; dans le dernier membre, chacun des deux signes moins est écrit au-dessus d'une tache qui couvre le signe d'abord posé. Dans la dernière équation, entre « $-$ » et $\frac{a^3(r^2-a^2)}{2}$, un coefficient d'abord écrit (une fraction) est couvert de hachures verticales serrées. Ici le 2 sous $u$ est bien formé, ce qui confirme la lecture des vues 23–29.}

On voit ici en effet qu'en faisant $r = a$ le second terme de l'equation disparoit, et qu'on a rigoureusement quel que soit $y'$, l'equation $\frac{u}{2} + \frac{a\,d u}{d r} = 0$. Donc puisque d'un autre coté la valeur de $\frac{u}{2} + \frac{a\,d u}{d r}$ n'est pas toujours nulle, il faut que le terme multiplié par $r^2 - a^2$ ne \struck{\uncertain{soit}} disparoisse pas toujours en faisant $r = a$.
\note{Le texte s'arrête sur cette phrase, aux deux tiers de la page ; le bas du feuillet est blanc. La vue 33 commence un nouveau paragraphe, « 7. ».}

\page{33}
\note{En haut à droite, « 8 » à l'encre. La page ouvre un nouveau paragraphe, qui suit la fin de phrase de la vue 31.}

7. Pour aller du plus simple au plus composé nous commencerons par examiner le cas où $y'$ est une quantité constante. Dans ce cas on n'a qu'à chercher l'integrale de \struck{\ill{}} $\rho^3\,d\pi'\,d\theta' \sin\theta'$, et pour cela nous pouvons supposer comme dans le n\textsuperscript{o}. 3, l'angle $\theta$ nul, ce qui donne $\mu = \cos\theta'$ et $d\mu = -\sin\theta'\,d\theta'$; de sorte que la differentielle dont il s'agit deviendra
\[ -\frac{d\pi'\,d\mu}{(r^2 - 2ar\mu + a^2)^{\frac{3}{2}}}. \]
Integrant d'abord par rapport à $\pi'$ on a
\[ -\frac{2\pi\,d\mu}{(r^2 - 2ar\mu + a^2)^{\frac{3}{2}}}; \]
integrant ensuite par rapport à $\mu$ on a
\[ \frac{-2\pi}{ar\sqrt{(r^2 - 2ar\mu + a^2)}}; \]
et completant depuis $\mu = 1$ jusqu'à $\mu = -1$, on aura l'integrale complette
\[ \frac{-2\pi}{ar(r+a)} + \frac{2\pi}{ar(r-a)} = \struck{\ill{}}\,\frac{4\pi}{r(r^2 - a^2)}. \]
Ainsi la valeur complette de $\int \rho^3 y'\,d\pi'\,d\theta' \sin\theta'$ lorsque $y'$ est constant, est $\struck{\ill{}}\,\frac{4\pi y'}{r(r^2 - a^2)}$ quelle que soit la valeur de $r$.
\note{Avant $\rho^3$, au début de l'intégrale, une lettre noyée sous une tache (un premier $\rho$ ?). « n\textsuperscript{o}. 3 » renvoie au paragraphe 3 (vues 23–25), où l'axe des angles est porté sur la ligne $r$. Le 2 du premier numérateur est empâté, et le $r$ du premier dénominateur porte une tache. Devant $\frac{4\pi}{r(r^2-a^2)}$ et devant $\frac{4\pi y'}{r(r^2-a^2)}$, une petite tache carrée, peut-être un signe effacé.}

Substituant cette valeur dans l'equation du numero precedent on aura
\[ \frac{u}{2} + \frac{r\,d u}{d r} = \struck{\ill{}} = -\frac{2\pi a^3 y'}{r} \]
Si maintenant on fait $r = a$, elle devient
\[ \frac{u}{2} + \frac{a\,d u}{d r} = -2a^2\pi y' \]
comme on l'a trouvée dans le n\textsuperscript{o}. 3, les quantités $y'$ et $y$ étant ici les mêmes.
\note{Le premier second membre de l'équation, où l'on devine $\pi$, $a$ et $y'$, est couvert de hachures et barré d'un double trait. « l'a trouvée » : ainsi, pour l'équation. Le renvoi au « n\textsuperscript{o}. 3 » vise la valeur $-2\pi a^2 y'$ de la vue 25, qui est en fait au paragraphe 3. La page s'achève sur cette phrase, au pied du feuillet.}

\page{35}
\note{En haut à droite, « 9 » à l'encre, et au-dessous, un peu à gauche, un petit signe en forme de c. La page ouvre un nouveau paragraphe, après la fin de la vue 33.}

8 Considerons maintenant le cas general dans lequel $y'$ est supposée une fonction des sinus et des cosinus des angles $\pi'$ et $\theta'$; on aura à integrer la quantité $\rho^3 y'\,d\pi'\,d\theta' \sin\theta'$; et pour cela nous transporterons aussi comme dans le numero 3, l'axe ou le pole des angles variables $\pi'$, $\theta'$ qui determinent la position de chaque molecule sur la sphere, dans la ligne $r$ menée du centre au point attiré; nommons $\varphi$ et $\psi$ ce que deviennent alors les angles $\pi'$ et $\theta'$, on aura egalement $d\varphi\,d\psi \sin\psi$ pour \struck{le} l'element $d\pi'\,d\theta' \sin\theta'$, et \struck{\ill{}} $\cos\psi$ pour la valeur de $\mu$; de sorte que la differentielle dont il s'agit deviendra $-\rho^3 y'\,d\varphi\,d\mu$, qu'il faudra integrer depuis $\varphi = 0$ jusqu'à $\varphi = 2\pi$ et depuis $\mu = 1$ jusqu'a $\mu = -1$.
\note{« 8 » : sans point. Les deux nouveaux angles : le premier est une boucle traversée d'une haste verticale, transcrite $\varphi$ ; le second, une croix dont la branche droite remonte en crochet, transcrite $\psi$ ; ni l'une ni l'autre forme n'est celle du $\theta$ ou du $\pi$ de la main. Le mot biffé avant « $\cos\psi$ » est rayé d'un trait épais qui le couvre entièrement.}

Mais comme $y'$ est supposée fonction des sinus et cosinus des angles $\pi'$ et $\theta'$ qui ont leur pole dans un axe fixe independant de la position du point attiré, ou de la ligne $r$, il faudra substituer dans $y'$ à la place des sinus et cosinus de $\pi'$ et $\theta'$ leurs valeurs en sinus et cosinus de $\varphi$ et $\psi$. Pour cela il n'y a qu'à considerer le triangle spherique qui a les angles $\theta$, $\theta'$ et $\psi$ pour ses trois cotés
\note{Entre « dans » et « un axe », de petites lettres serrées au-dessus de la ligne, non lues. Le premier $\theta$ de la dernière ligne porte au-dessus un petit trait empâté, qui n'est pas le prime du $\theta'$ suivant. La page s'arrête sur « cotés », sans ponctuation, aux trois quarts du feuillet ; la phrase continue en tête de la vue 37.}

\page{37}
\note{En haut à droite, « 10 » à l'encre. La page continue la phrase laissée sur « cotés » à la fin de la vue 35.}

et dans lequel $\pi - \pi'$ est l'angle opposé au coté $\psi$, et $\varphi$ est l'angle opposé au coté $\theta'$; et les formules connues de la trigonometrie spherique donneront
\[ \begin{array}{l} \cos\theta' = \cos\theta\cos\psi + \sin\theta\sin\psi\cos\varphi \\ \sin\theta'\sin(\pi - \pi') = \sin\psi\sin\varphi \\ \sin\theta'\cos(\pi - \pi') = \sin\theta\cos\psi + \cos\theta\sin\psi\cos\varphi \end{array} \]

\marginal{I 65}
\note{Dans la marge de droite, en regard de la ligne qui suit, « I 65 » encadré d'un trait fin, d'une plume plus mince que celle du texte ; au début de la même ligne, un crochet ouvrant « [ » devant « Ainsi ». Ce que marquent ce crochet et ce chiffre (un renvoi, une marque de copie ou d'impression ?), la page ne le dit pas ; aucun crochet fermant n'a été vu dans le lot.}

Ainsi la formule $\rho^3 y'\,d\varphi\,d\mu$ sera integrable toutes les fois que $y'$ sera une fonction rationnelle et entiere de $\cos\theta'$, $\sin\theta'\sin\pi'$, $\sin\theta'\cos\pi'$, parceque l'integration relative à $\varphi$ fera disparoitre toutes les puissances impaires de $\sin\varphi$. C'est ainsi que d'Alembert a calculé l'attraction des spheroides peu differens de la sphere; mais pour notre objet il suffit d'avoir reduit l'integration de la formule \struck{\ill{}}$\int \rho^3 y'\,d\pi'\,d\theta' \sin\theta'$ à celle de la formule $-\rho^3 y'\,d\varphi\,d\mu$. \struck{\uncertain{depuis} \ill{}}
\note{« parceque » : en un mot. Un petit signe devant « $\sin\theta'\sin\pi'$ ». Devant l'intégrale, un trait court et épais, biffé. Après « $d\varphi\,d\mu$. », deux mots rayés d'un trait continu ; le premier se lit « depuis » avec doute.}

9 Commençons par l'integration relative à $\mu$; la differentielle $y'\rho^3d\mu$ integrée par parties donne
\[ y'\int\rho^3d\mu - \int \frac{d y'}{d\mu}\left(\int\rho^3d\mu\right)d\mu. \]
Or $\rho = \frac{1}{\sqrt{(r^2 - 2ar\mu + a^2)}}$ (n\textsuperscript{o}. 4) donc $\int\rho^3d\mu = \frac{\rho}{ar}$; de sorte que l'integrale dont il s'agit deviendra
\[ \frac{y'\rho}{ar} - \int \frac{d y'}{d\mu} \times \frac{\rho\,d\mu}{ar}. \]
Comme l'integration doit s'etendre depuis $\mu = 1$ jusqu'a $\mu = -1$, et qu'a la premiere de ces limites $y'$ devient $y$, si on nomme $Y$ ce que $y'$ devient à la seconde
\note{« 9 » : sans point. « n\textsuperscript{o}. 4 » renvoie au paragraphe 4 (vue 25), où $\rho$ est défini. « $Y$ » : la lettre a la forme d'un 4 bouclé, la même que l'angle $\psi$ de la vue 35 ; c'est la capitale $Y$ de la vue 39, où la même quantité est écrite aussi en capitale barrée. La page s'arrête là, sans ponctuation, aux trois quarts du feuillet ; la phrase continue en tête de la vue 39.}

\page{39}
\note{En haut à droite, « 11 » à l'encre. La page continue la phrase laissée à la fin de la vue 37 (« à la seconde »).}

limite, la valeur complette du terme $\frac{y'\rho}{ar}$ hors du signe sera
\[ \frac{Y}{ar(r+a)} - \frac{y}{ar(r-a)}. \]
Or $y$ et $Y$ repondant à $\psi = 0$ et $\psi = \pi$ à cause de $\mu = \cos\psi$, auquel cas $\pi'$ et $\theta'$ deviennent $\pi$ et $\theta$, il est visible que ces quantités seront independantes de $\varphi$; ainsi l'integration relative à $\varphi$ depuis $\varphi = 0$ jusqu'à $\varphi = 2\pi$, executée sur les memes termes donnera
\[ \frac{2\pi Y}{ar(r+a)} - \frac{2\pi y}{ar(r-a)}. \]
On pourra donc substituer dans l'equation \ill{} du n\textsuperscript{o}. 6, à la place de $\int \rho^3 y'\,d\pi'\,d\theta' \sin\theta'$ la quantité
\[ \frac{2\pi y}{ar(r-a)} - \frac{2\pi Y}{ar(r+a)} + \int \frac{d y'}{d\mu} \times \frac{\rho\,d\mu\,d\varphi}{ar} \]
ce qui la changera en celle-ci
\[ \frac{u}{2} + \frac{r\,d u}{d r} = -\frac{a^2\pi(r+a)y}{r} + \frac{a^2\pi Y(r-a)}{r} - \frac{a^3(r^2 - a^2)}{2ar} \int \frac{d y'}{d\mu}\,\rho\,d\mu\,d\varphi; \]
dans laquelle il faut faire maintenant $r = a$ ce qui la reduit à celle-ci
\[ \frac{u}{2} + \frac{a\,d u}{d r} = -2a^2\pi y. \]
comme dans le cas où $y$ est constant.
\note{$Y$ (la valeur de $y'$ à la seconde limite) est une capitale barrée d'un trait horizontal dans les deux premières fractions ; dans l'avant-dernière équation elle prend la forme d'un 4 bouclé, la même qu'à la fin de la vue 37 et que l'angle $\psi$ ; les deux lettres se confondent sur la page et ne se séparent que par le sens. Le mot entre « l'equation » et « du n\textsuperscript{o}. 6 » n'a pas été lu (« avec ce » ? « ancienne » ?). Le n\textsuperscript{o}. 6 est celui de la vue 31, où l'équation $\frac{u}{2} + \frac{r\,d u}{d r} = -\frac{a^3(r^2-a^2)}{2}\int\rho^3 y'\ldots$ est établie.}

Il est bon de remarquer qu'une nouvelle integration par parties executée sur la differentielle $\frac{d y'}{d\mu}\,\rho\,d\mu$ \struck{\ill{}} n'ajouteroit \struck{\ill{}} rien
\note{Deux mots courts rayés d'un trait épais, avant et après « n'ajouteroit ». La page s'arrête sur « rien », sans ponctuation, aux deux tiers du feuillet ; le bas est blanc, et l'on y voit par transparence, à l'envers, l'écriture du verso (vue 40). Ce lot ne dit pas si la phrase continue à la vue 41.}

\page{40}
\note{Verso du feuillet de la vue 39. Le haut de la vue ne montre que le texte du recto, par transparence. Au bas, tête-bêche (la vue retournée de 180° se lit normalement), trois lignes de la main du texte, coupées de quatre longs traits obliques qui les annulent ; à leur tête, dans le coin qui devient alors le coin supérieur droit, un « 11 » à l'encre traversé d'un trait oblique. C'est un premier départ de la page 11 : ses mots sont ceux du début de la vue 39. Il est transcrit ici tel qu'on le lit la page retournée.}

\struck{limite, la valeur complette du terme $\frac{y'\rho}{ar}$ hors du}

\struck{\ill{} \uncertain{signe sera} $\frac{Y(r+a)}{ar} - \frac{y(r\uncertain{+}a)}{ar}$}

\struck{signe sera $\frac{Y}{ar(r+a)} - \frac{y}{ar(r-a)}$; ainsi}
\note{La deuxième ligne est surchargée : ses premiers mots sont noyés d'encre, et ses deux fractions, récrites l'une sur l'autre, sont reprises à la troisième ligne sous une autre forme, celle que garde la vue 39. Le signe dans le second binôme de la deuxième ligne est douteux. Le dernier mot, « ainsi », porte en plus un trait qui lui est propre.}

\note{À droite, sous ces lignes (la page retournée), un croquis à la plume : deux courbes presque parallèles en arc, qui se rejoignent en pointe à une extrémité, coupées vers leur milieu par un petit trait transversal ; aucune lettre. Ce que le croquis représente, la page ne le dit pas.}

\end{document}
